\par
\SourceLeaf{378}%
Noch ein paar Bemerkungen glauben wir beifügen zu müssen. Wir finden bei sämmtlichen Messungen, dass man beim Anfange jeder Serie immer den Index auf 0 zurückbrachte, ein Verfahren, welches wir nicht billigen können, weil dadurch, wie sehr man auch die Anzahl der Serien vervielfältigt, immer derselbe vom Theilungsfehler abhängige constante Fehler im Resultate zurückbleiben muss. --- Bei den Messungen der Piemontesischen Officiere ist jedesmal der Zustand der Luft angezeigt. Unter 414 Messungsreihen zählen wir 320, wo Windstille, und 94, wo Wind angezeigt ist: ein so günstiges Verhältniss hätte man an so hochliegenden Standpunkten (die Höhe des höchsten über der Meeresfläche beträgt 3534 Meter) kaum erwartet.

Der zweite Band enthält in zehn Abschnitten die Arbeiten der Astronomen. In den beiden ersten Abschnitten finden wir die auf die Längengradmessung Beziehung habenden Bestimmungen von Längenunterschieden durch Pulversignale. Die ersten Versuche dieser Art wurden im September 1821 gemacht; die Pulversignale wurden auf der Rocca Melone gegeben, und auf der 170,000 Meter entfernten Sternwarte von Mailand und auf dem nahen Mont Cenis beobachtet. Für die Zeitbestimmung an letzterm Platze war in dem Garten des Hospizes eine kleine Sternwarte errichtet und ein Mittagsfernrohr von \textsc{Fortin} darin aufgestellt, welches jedoch nicht von ausgezeichneter Güte gewesen zu sein scheint, wie in Beziehung auf die Zapfen, einen wesentlichen Theil, ausdrücklich bemerkt wird. Die Rocca Melone war hier nicht sichtbar; man musste sich, um die Signale zu sehen, an eine etwas entfernte Stelle begeben, wohin man die Zeit mit einem Chronometer von \textsc{Earnshaw} übertrug. Auch die Beobachtungen am Mittagsfernrohre wurden meistens an diesem Chronometer notirt, aber nicht vom Beobachter selbst, sondern nach einem von diesem gegebenen Zeichen, durch einen Gehülfen. Alle diese Umstände vereinigen sich freilich, das Zutrauen zu der. Genauigkeit des Endresultats zu verringern, wenn gleich die drei partiellen Resultate von den drei Beobachtungstagen sehr gut übereinstimmen. Es kommt dazu, dass man hier die Zeitbestimmung aus Sternen, in Mailand aus Sonnendurchgängen erhielt, und endlich, dass, wie es scheint, die Rechtwinkligkeit der optischen Axe des \textsc{Fortin}schen Mittagsfernrohrs zu dessen Drehungsaxe gar nicht berichtigt wurde, wenigstens wird dieses wichtigen Umstandes bei diesen Beobachtungen gar nicht erwähnt.

Bei den Operationen ähnlicher Art im Jahr 1822 ging man in jeder Bezie%
